\documentclass[11pt]{article}

\usepackage[margin=1in]{geometry}
\usepackage{graphicx}
\usepackage{amsmath}
\usepackage{xr}
\externaldocument{main}
\usepackage[colorlinks=true,allcolors=blue]{hyperref}

% Numbers the manuscript prints.
%
% Two regions. Everything between the markers below is regenerated from the
% committed result files and is never hand-edited. Anything outside them is a
% fixed convention this repository did not measure -- a value quoted from a
% cited paper -- and is hand-edited and reviewed like any other source.
%
% Each generated block declares the population its macros are measured over: the
% artifact the number came from, and the unit noun a count of that population
% takes. Two population names resolving to one artifact is a test failure.

%% POPULATION cited | artifact: npj Comput. Mater. 11, 300 (2025) | unit: rows
% Quoted from Kim & Lee's article, not measured here.
\newcommand{\npjSystems}{43}

%% POPULATION citedRelease | artifact: AI Hub dataset 71339 and Matbench v0.1 | unit: releases
% EMPD's release number, as its data description prints it, and the date on
% which the local copy was downloaded, which the archive's directory records.
% Matbench's version, as Dunn et al. name the suite that holds the task.
\newcommand{\empdRelease}{1.0}
\newcommand{\empdRetrieved}{23 December 2025}
\newcommand{\exptGapVersion}{0.1}

% ---- BEGIN GENERATED ----

%% POPULATION npjDataset | artifact: data/npj/csv/Train.csv+data/npj/csv/Val.csv | unit: rows
\newcommand{\npjDatasetRows}{71{,}633}

%% POPULATION npjTrain | artifact: data/npj/csv/Train.csv | unit: rows
\newcommand{\npjModelTrainRmse}{3.95}
\newcommand{\npjPaperTrainRmse}{3.94}

%% POPULATION npjVal | artifact: data/npj/csv/Val.csv | unit: rows
\newcommand{\npjModelRSq}{0.9946}
\newcommand{\npjModelRmse}{5.87}
\newcommand{\npjPaperRmse}{5.86}
\newcommand{\npjFeatureCount}{nine} % artifact: data/npj/csv/Val.csv | unit: features
\newcommand{\npjCompositionFeatureCount}{six} % artifact: data/npj/csv/Val.csv | unit: descriptors
\newcommand{\npjLookupRSq}{0.9924}
\newcommand{\npjMargin}{+0.0022}
\newcommand{\npjLookupRmse}{6.94}

%% POPULATION empdRes | artifact: outputs/aihub/draws.json:resistivity | unit: rows
\newcommand{\empdResBlockedHeldOutMin}{9\%}
\newcommand{\empdResBlockedHeldOutMax}{20\%}
\newcommand{\empdResRows}{58{,}937}
\newcommand{\empdResSystems}{25} % artifact: outputs/aihub/draws.json:resistivity | unit: systems
\newcommand{\empdResLookup}{0.887}
\newcommand{\empdResModel}{0.946}
\newcommand{\empdResMargin}{+0.059}
\newcommand{\empdResMarginCI}{[+0.057, +0.061]}
\newcommand{\empdResBlockedLookup}{$-$0.1}
\newcommand{\empdResBlockedLookupCI}{[$-$0.6, 0.4]}
\newcommand{\empdResBlocked}{0.1}
\newcommand{\empdResBlockedCI}{[$-$0.2, 0.4]}
\newcommand{\empdResExcess}{+0.8}
\newcommand{\empdResExcessCI}{[+0.4, +1.1]}
\newcommand{\empdResBlockedMin}{$-$0.936}
\newcommand{\empdResBlockedMax}{+0.839}
\newcommand{\empdResMarginMin}{+0.053}
\newcommand{\empdResMarginMax}{+0.064}
\newcommand{\empdResBlockedNegative}{three} % artifact: outputs/aihub/draws.json:resistivity | unit: partitions
\newcommand{\empdResBlockedPositive}{seven} % artifact: outputs/aihub/draws.json:resistivity | unit: partitions
\newcommand{\empdResValidationOnlySystems}{one} % artifact: outputs/aihub/draws.json:resistivity | unit: systems

%% POPULATION empdHard | artifact: outputs/aihub/draws.json:hardness | unit: rows
\newcommand{\empdHardBlockedHeldOutMin}{6\%}
\newcommand{\empdHardBlockedHeldOutMax}{20\%}
\newcommand{\empdHardRows}{901}
\newcommand{\empdHardSystems}{22} % artifact: outputs/aihub/draws.json:hardness | unit: systems
\newcommand{\empdHardLookup}{0.873}
\newcommand{\empdHardModel}{0.921}
\newcommand{\empdHardMargin}{+0.048}
\newcommand{\empdHardMarginCI}{[+0.035, +0.063]}
\newcommand{\empdHardMarginMin}{+0.017}
\newcommand{\empdHardBlockedLookup}{$-$4}
\newcommand{\empdHardBlockedLookupCI}{[$-$10, 0]}
\newcommand{\empdHardBlocked}{$-$6}
\newcommand{\empdHardBlockedCI}{[$-$12, $-$1]}
\newcommand{\empdHardExcess}{+7}
\newcommand{\empdHardExcessCI}{[+2, +13]}
\newcommand{\empdHardBlockedMin}{$-$23.541}
\newcommand{\empdHardBlockedMax}{+0.730}
\newcommand{\empdHardBlockedMedian}{$-$1.354}
\newcommand{\empdHardBlockedNegative}{nine} % artifact: outputs/aihub/draws.json:hardness | unit: partitions
\newcommand{\empdHardMarginMax}{+0.103}
\newcommand{\empdHardValidationOnlySystems}{four} % artifact: outputs/aihub/draws.json:hardness | unit: systems

%% POPULATION empdAmorph | artifact: outputs/aihub/draws.json:amorphous | unit: rows
\newcommand{\empdAmorphousRows}{945}

%% POPULATION npjShared | artifact: outputs/aihub/structure.json:npj_shared_measurements | unit: rows
\newcommand{\npjSharedRows}{2{,}405}
\newcommand{\npjSharedOtherRows}{157}
\newcommand{\npjSharedFactor}{0.45326}
\newcommand{\npjFactorGridLow}{0.38}
\newcommand{\npjFactorGridHigh}{0.52}
\newcommand{\npjFactorGridStep}{0.00001}
\newcommand{\npjSharedMinDecimals}{six} % artifact: outputs/aihub/structure.json:npj_shared_measurements | unit: decimals
\newcommand{\npjSharedThicknessRatio}{1.000}
\newcommand{\npjSharedThicknessRows}{2{,}289}

%% POPULATION htemRes | artifact: outputs/htem/split.json:log10_resistivity | unit: rows
\newcommand{\htemResRows}{8{,}628}
\newcommand{\htemResLibraries}{257} % artifact: outputs/htem/split.json:log10_resistivity | unit: libraries
\newcommand{\htemResRandomLookup}{0.84}
\newcommand{\htemResRandomModel}{0.90}
\newcommand{\htemResBlockedLookup}{0.14}
\newcommand{\htemResBlockedLookupCI}{[$-$0.06, 0.31]}
\newcommand{\htemResBlockedModel}{0.36}
\newcommand{\htemResBlockedModelCI}{[0.20, 0.48]}
\newcommand{\htemResMargin}{+0.057}
\newcommand{\htemResMarginCI}{[+0.042, +0.075]}
\newcommand{\htemResExcess}{+0.48}
\newcommand{\htemResExcessCI}{[+0.37, +0.62]}
\newcommand{\htemResBeyondHalf}{+0.28}
\newcommand{\htemResBeyondHalfCI}{[+0.11, +0.48]}
\newcommand{\htemResFallExcess}{+0.16}
\newcommand{\htemResFallExcessCI}{[+0.04, +0.29]}
\newcommand{\htemResFillFallLo}{+0.037}
\newcommand{\htemResMarginMin}{+0.051}
\newcommand{\htemResMarginMax}{+0.064}

%% POPULATION htemGap | artifact: outputs/htem/split.json:direct_band_gap | unit: rows
\newcommand{\htemGapRows}{5{,}264}
\newcommand{\htemGapLibraries}{156} % artifact: outputs/htem/split.json:direct_band_gap | unit: libraries
\newcommand{\htemGapRandomLookup}{0.66}
\newcommand{\htemGapRandomLookupCI}{[0.56, 0.74]}
\newcommand{\htemGapRandomLookupLo}{0.559}
\newcommand{\htemGapRandomModel}{0.78}
\newcommand{\htemGapBlockedLookup}{0.2}
\newcommand{\htemGapBlockedLookupCI}{[$-$0.1, 0.4]}
\newcommand{\htemGapBlockedModel}{0.47}
\newcommand{\htemGapBlockedModelCI}{[0.33, 0.58]}
\newcommand{\htemGapMargin}{+0.12}
\newcommand{\htemGapMarginCI}{[+0.09, +0.16]}
\newcommand{\htemGapExcess}{+0.19}
\newcommand{\htemGapExcessCI}{[+0.10, +0.29]}
\newcommand{\htemGapBeyondHalf}{+0.16}
\newcommand{\htemGapBeyondHalfCI}{[$-$0.01, +0.37]}
\newcommand{\htemGapFallExcess}{+0.19}
\newcommand{\htemGapFallExcessCI}{[+0.06, +0.33]}
\newcommand{\htemGapFillFallLo}{+0.018}
\newcommand{\htemGapMarginMin}{+0.114}
\newcommand{\htemGapMarginMax}{+0.130}
\newcommand{\htemGapFirstMargin}{+0.13}
\newcommand{\htemGapFirstMarginCI}{[+0.09, +0.17]}

%% POPULATION htemResAblation | artifact: outputs/htem/ablation.json:log10_resistivity | unit: rows
\newcommand{\htemResRandomTuned}{0.90}
\newcommand{\htemResBlockedTuned}{0.39}
\newcommand{\htemResExcessTuned}{+0.45}
\newcommand{\htemResExcessTunedCI}{[+0.35, +0.56]}
\newcommand{\htemResRandomBoosted}{0.90}
\newcommand{\htemResBlockedBoosted}{0.33}
\newcommand{\htemResExcessBoosted}{+0.51}
\newcommand{\htemResExcessBoostedCI}{[+0.39, +0.64]}

%% POPULATION htemGapAblation | artifact: outputs/htem/ablation.json:direct_band_gap | unit: rows
\newcommand{\htemGapRandomTuned}{0.79}
\newcommand{\htemGapBlockedTuned}{0.55}
\newcommand{\htemGapExcessTuned}{+0.11}
\newcommand{\htemGapExcessTunedCI}{[+0.03, +0.20]}
\newcommand{\htemGapRandomBoosted}{0.79}
\newcommand{\htemGapBlockedBoosted}{0.48}
\newcommand{\htemGapExcessBoosted}{+0.18}
\newcommand{\htemGapExcessBoostedCI}{[+0.09, +0.27]}

%% POPULATION exptGap | artifact: outputs/expt_gap/control.json | unit: rows
\newcommand{\exptGapRows}{4{,}604}
\newcommand{\exptGapLookup}{0.26}
\newcommand{\exptGapLookupCI}{[0.19, 0.32]}
\newcommand{\exptGapModel}{0.48}
\newcommand{\exptGapMargin}{+0.23}
\newcommand{\exptGapMarginCI}{[+0.17, +0.29]}
\newcommand{\exptGapMarginMin}{+0.193}
\newcommand{\exptGapMarginMax}{+0.267}
\newcommand{\exptGapFirstMargin}{+0.19}
\newcommand{\exptGapFirstMarginCI}{[+0.13, +0.26]}
\newcommand{\exptGapZeroShare}{53\%}
\newcommand{\exptGapNonzeroRows}{2{,}154} % artifact: outputs/expt_gap/control.json | unit: compounds
\newcommand{\exptGapNonzeroLookup}{0.51}
\newcommand{\exptGapNonzeroLookupCI}{[0.45, 0.56]}
\newcommand{\exptGapNonzeroLookupHi}{0.562}
\newcommand{\exptGapNonzeroMargin}{+0.20}
\newcommand{\exptGapNonzeroMarginCI}{[+0.14, +0.26]}

%% POPULATION programs | artifact: reproduce/make_parameters.py | unit: sources
\newcommand{\programCount}{two}

%% POPULATION design | artifact: reproduce/ | unit: settings
\newcommand{\mlpFirstLayer}{128} % artifact: reproduce/ | unit: units
\newcommand{\mlpSecondLayer}{64} % artifact: reproduce/ | unit: units
\newcommand{\mlpDropout}{0.3}
\newcommand{\mlpEpochs}{200} % artifact: reproduce/ | unit: epochs
\newcommand{\mlpPatience}{30} % artifact: reproduce/ | unit: epochs
\newcommand{\mlpBatchSize}{32} % artifact: reproduce/ | unit: rows
\newcommand{\mlpLearningRate}{$10^{-3}$}
\newcommand{\mlpWeightDecay}{$10^{-5}$}
\newcommand{\mlpLrFactor}{0.5}
\newcommand{\mlpLrPatience}{10} % artifact: reproduce/ | unit: epochs
\newcommand{\mlpMaxGradNorm}{1.0}
\newcommand{\randomPairs}{five} % artifact: outputs/ (the pairs that the result files score) | unit: pairs
\newcommand{\libraryDatasets}{three} % artifact: outputs/ (the pairs that the result files score) | unit: datasets
\newcommand{\largestRandomMargin}{+0.12} % artifact: outputs/ (the pairs that the result files score) | unit: settings
\newcommand{\largestRandomMarginPair}{HTEM-DB band gap} % artifact: outputs/ (the pairs that the result files score) | unit: settings
\newcommand{\blockedPairs}{four} % artifact: outputs/ (the pairs that the result files score) | unit: pairs
\newcommand{\blockedDatasets}{two} % artifact: outputs/ (the pairs that the result files score) | unit: datasets
\newcommand{\splitDraws}{ten} % artifact: reproduce/ | unit: partitions
\newcommand{\splitInits}{three} % artifact: reproduce/ | unit: times
\newcommand{\heldOutShare}{15\%}
\newcommand{\empdTrainShare}{70\%}
\newcommand{\empdStopShare}{15\%}
\newcommand{\npjCommit}{8c2f808}
\newcommand{\htemFolds}{five} % artifact: reproduce/ | unit: folds
\newcommand{\htemTrees}{200} % artifact: reproduce/ | unit: trees
\newcommand{\htemNoTemperature}{$-$1}
\newcommand{\htemNoTemperatureAlternatives}{0, $-$100 or $-$1000}
\newcommand{\htemNoTemperatureShift}{0.0049} % artifact: outputs/htem/split.json | unit: settings
\newcommand{\htemNoTemperatureLookupShift}{0.069} % artifact: outputs/htem/split.json | unit: settings
\newcommand{\empdElementCount}{15} % artifact: reproduce/ | unit: elements
\newcommand{\empdDescriptorCount}{six} % artifact: reproduce/ | unit: descriptors
\newcommand{\htemTunedTrees}{400} % artifact: reproduce/ | unit: trees
\newcommand{\htemSearchBudget}{25} % artifact: reproduce/ | unit: configurations
\newcommand{\htemInnerFolds}{four} % artifact: reproduce/ | unit: folds
\newcommand{\htemBoostRounds}{1{,}000} % artifact: reproduce/ | unit: rounds
\newcommand{\htemForestFeatureShares}{0.05, 0.15, 0.3, 0.5 or 1.0}
\newcommand{\htemForestMinLeaf}{1, 3, 10, 30 or 100}
\newcommand{\htemBoostRates}{0.01, 0.03, 0.06, 0.12 or 0.25}
\newcommand{\htemBoostLeaves}{3, 7, 15, 63 or 255}
\newcommand{\htemBoostPenalty}{1.0}
\newcommand{\bootstrapReplicates}{2{,}000} % artifact: reproduce/ | unit: resamples
\newcommand{\bootstrapLevel}{95}

%% POPULATION software | artifact: requirements.txt | unit: versions
\newcommand{\numpyVersion}{1.26.4}
\newcommand{\scipyVersion}{1.13.1}
\newcommand{\sklearnVersion}{1.6.1}
\newcommand{\torchVersion}{2.8.0}

% ---- END GENERATED ----


\usepackage{titlesec}
\renewcommand{\thesection}{S\arabic{section}}
\titleformat{\section}{\normalfont\large\bfseries}{Section~\thesection.}{0.5em}{}
\renewcommand{\figurename}{Fig.}
\renewcommand{\thefigure}{S\arabic{figure}}
\renewcommand{\thetable}{S\arabic{table}}

\title{Supplementary Information for\\
       \emph{Two reporting checks for machine learning accuracy on
       combinatorial materials libraries}}
\author{Sungwoo Kang\\
  Department of Electrical and Computer Engineering, Korea University,\\
  Seoul, Republic of Korea\\[2ex]
  Jong-Kook Kim\\
  School of Electrical Engineering, Korea University,\\
  Seoul, Republic of Korea}
\date{}

\begin{document}
\maketitle

\section{Spread of scores across the partitions}
\label{sec:esi-draw-spread}

The range of a margin across the \splitDraws\ random partitions shows how far
one partition's margin can lie from the mean over them, the value that the main
text reports on EMPD, HTEM-DB and Matbench. On EMPD hardness, that
mean is \empdHardMargin\ and the partitions give margins from
\empdHardMarginMin\ to \empdHardMarginMax, so one partition's margin can be
more than twice the mean. The margin varies less across the random partitions
on the other pairs of EMPD and HTEM-DB and on Matbench, from \empdResMarginMin\ to \empdResMarginMax\ on
EMPD resistivity, from \htemResMarginMin\ to \htemResMarginMax\ on HTEM-DB
resistivity, from \htemGapMarginMin\ to \htemGapMarginMax\ on HTEM-DB band gap
and from \exptGapMarginMin\ to \exptGapMarginMax\ on Matbench.

On Matbench, the lookup scores \exptGapLookup\
\exptGapLookupCI\ (bracketed ranges are \bootstrapLevel\% percentile bootstrap
intervals), below the \htemGapRandomLookup\ \htemGapRandomLookupCI\ on HTEM-DB band
gap. Because \exptGapZeroShare\ of Matbench's band gaps are zero and no HTEM-DB
band gap is, we also draw \splitDraws\ random partitions of the
\exptGapNonzeroRows\ Matbench compounds whose gap exceeds zero, each from the seed
of a random partition of the full task, and score the lookup and the forest on
them. Under these partitions, however, the lookup scores \exptGapNonzeroLookup\
\exptGapNonzeroLookupCI. The upper end of that interval, \exptGapNonzeroLookupHi,
exceeds the lower end of HTEM-DB band gap's interval, \htemGapRandomLookupLo, so the
zero gaps, and not only the absence of libraries, can lower the lookup's score on
Matbench.

The Matbench margin's interval lies
above that of HTEM-DB band gap, the pair that Matbench matches
(Table~\ref{tab:main}a of the main text). However, we first compared the two margins on a
single random partition, the first of the \splitDraws. We moved the comparison
to the means over the \splitDraws\ partitions only after the intervals of the two margins overlapped on
that partition, so on Matbench the wider margin is not established. On that
partition, the Matbench margin is \exptGapFirstMargin\ \exptGapFirstMarginCI\
and that of HTEM-DB band gap is \htemGapFirstMargin\ \htemGapFirstMarginCI. The first partition gives Matbench its smallest margin and HTEM-DB band gap
one of its two largest, which explains why the two intervals overlap on it but not
over the means. On the Matbench compounds whose gap exceeds zero, the margin is \exptGapNonzeroMargin\ \exptGapNonzeroMarginCI, an interval that overlaps HTEM-DB band gap's interval even over the means, so
on those compounds the wider margin is not established either.

On EMPD, the model's blocked $R^2$ varies so widely across the \splitDraws\
blocked partitions that its mean describes no typical partition, although the
mean still compares the model with a constant prediction of the test rows' mean
target. The blocked $R^2$ ranges from \empdResBlockedMin\ to
\empdResBlockedMax\ on resistivity and from \empdHardBlockedMin\ to
\empdHardBlockedMax\ on hardness, partly because each partition's $R^2$ is
computed over a different set of test rows, relative to the variance of their
targets. Because of this spread, the model's mean on hardness in
Table~\ref{tab:main}b of the main text, \empdHardBlocked, lies far below its
median of \empdHardBlockedMedian, so the mean depends mostly on the partitions
with the lowest $R^2$. The model's blocked $R^2$ on hardness lies below zero, the
$R^2$ of a constant prediction of the test rows' mean, under
\empdHardBlockedNegative\ of the \splitDraws\ partitions, so under those
partitions the model predicts the hardness of the alloy systems excluded from
training less accurately than a constant prediction of the test rows' mean
hardness. The model's blocked $R^2$ on resistivity lies below zero under
\empdResBlockedNegative\ partitions and above zero under the other
\empdResBlockedPositive.

\section{Tuned models on the HTEM-DB blocked partition}
\label{sec:esi-posthoc}

To find whether the positive excess on HTEM-DB depends on the untuned forest, we
replace it with two models tuned inside their training folds, a random forest and
a gradient-boosting machine, both from scikit-learn. Each tuned model splits the
rows of its training folds into \htemInnerFolds\ inner folds and chooses among
\htemSearchBudget\ configurations by their scores on one of the inner folds. The
gradient-boosting machine also chooses on that inner fold its number of boosting
rounds, at most \htemBoostRounds. Unlike the untuned forest, which keeps its
settings under the blocked partition, each tuned model chooses its configuration
again under that partition, on inner folds that contain whole libraries. The
tuned forest averages \htemTunedTrees\ decision trees. Its configurations combine
a share of \htemForestFeatureShares\ of the features considered at each split with
a minimum of \htemForestMinLeaf\ rows per leaf. The configurations of the
gradient-boosting machine combine a learning rate of \htemBoostRates\ with at most
\htemBoostLeaves\ leaves per tree. Every configuration of the gradient-boosting
machine sets the $L_2$ regularization to \htemBoostPenalty. The choice of boosting
rounds and the grids correct two faults of an earlier version, found after that
version had reported its excess, so these corrections are post hoc. Under the
blocked partition, each tuned model chooses in some folds a configuration on an
edge of its grid, so a setting beyond the grid could change its excess.

The HTEM-DB excess is positive on both properties under either tuned model.
Every interval for that excess excludes zero, including that of the band-gap
excess, which the tuned forest lowers from \htemGapExcess\ under the untuned
forest to \htemGapExcessTuned\ (Table~\ref{tab:ablation}).

\begin{table}[htbp] \small \caption{Three models' scores under the HTEM-DB
random and blocked partitions, on the main text's rows and folds. Bracketed ranges are
\bootstrapLevel\% percentile bootstrap intervals over \bootstrapReplicates\
resamples of whole libraries.}
\label{tab:ablation} % Generated from the committed result files; not hand-edited.
\begin{tabular*}{\textwidth}{@{\extracolsep{\fill}}llrrl}
\hline
property & model & random $R^2$ & blocked $R^2$ & excess \\
\hline
resistivity & \htemTrees-tree forest, untuned & \htemResRandomModel & \htemResBlockedModel & \htemResExcess~\htemResExcessCI \\
 & \htemTunedTrees-tree forest, tuned & \htemResRandomTuned & \htemResBlockedTuned & \htemResExcessTuned~\htemResExcessTunedCI \\
 & gradient boosting, tuned & \htemResRandomBoosted & \htemResBlockedBoosted & \htemResExcessBoosted~\htemResExcessBoostedCI \\
\hline
band gap & \htemTrees-tree forest, untuned & \htemGapRandomModel & \htemGapBlockedModel & \htemGapExcess~\htemGapExcessCI \\
 & \htemTunedTrees-tree forest, tuned & \htemGapRandomTuned & \htemGapBlockedTuned & \htemGapExcessTuned~\htemGapExcessTunedCI \\
 & gradient boosting, tuned & \htemGapRandomBoosted & \htemGapBlockedBoosted & \htemGapExcessBoosted~\htemGapExcessBoostedCI \\
\hline
\end{tabular*}
 \end{table}

\section{Settings of the protocol}
\label{sec:esi-settings}

Each check compares scores within one dataset, so the checks do not need the
same model on every dataset. On Kim and Lee's dataset, we recompute their network's validation $R^2$ from their training and
validation files rather than quoting it, so that the network and the lookup are
scored on identical rows. On
those rows, the network's RMSE is \npjModelRmse~$\mu\Omega\cdot$cm and the
lookup's \npjLookupRmse~$\mu\Omega\cdot$cm, so, near an $R^2$ of one, RMSE
distinguishes the network from the lookup more clearly than $R^2$ does.

No divisor in a range around \npjSharedFactor, from \npjFactorGridLow\ to
\npjFactorGridHigh\ in steps of \npjFactorGridStep, reproduces more than
\npjSharedOtherRows\ of Kim and Lee's rows, apart from \npjSharedFactor\ itself,
so \npjSharedFactor\ is the conversion factor that relates Kim and Lee's
resistivity to EMPD's resistance. EMPD's raw files record a thickness at the
positions of \npjSharedThicknessRows\ of the \npjSharedRows\ rows that it
reproduces, where the median ratio of Kim and Lee's thickness to EMPD's is
\npjSharedThicknessRatio.

EMPD's release splits each table in row order, so a model trained on the
training part is tested partly on alloy systems excluded from that training,
because the validation part alone contains \empdResValidationOnlySystems\ of the
resistivity alloy systems and \empdHardValidationOnlySystems\ of the hardness
systems. That split is therefore not a random partition of the rows.

On HTEM-DB, any value below every recorded deposition temperature would serve
as the value of a row that records none, because replacing \htemNoTemperature\ by
\htemNoTemperatureAlternatives\ changes each HTEM-DB margin and excess by at most
\htemNoTemperatureShift\ in $R^2$ and each blocked-partition lookup score by at
most \htemNoTemperatureLookupShift\ in $R^2$. On both properties, the lookup's
$R^2$ still falls further than the model's under the blocked partition whichever
of these values replaces \htemNoTemperature. Under each of these values, the
interval of the lookup's fall minus the model's has a lower endpoint of at least
\htemResFillFallLo\ on resistivity and \htemGapFillFallLo\ on band gap.

\end{document}
